% !TEX root = sub-main.tex
\documentclass[main.tex]{subfiles}

\begin{document}
    \begin{frame}
        \frametitle{What Can Explain JP Morgan's XVA Loss?}
        
        \begin{figure}
            \centering
            \resizebox{\textwidth}{!}{%
            \begin{tikzpicture}
                \colorlet{funding-color}{wtf-blue}
\colorlet{rate-color}{wtf-red}

\tikzmath{
    \arrowyoffset = 7;
    \arrowcenteryoffset = 7;
    \toparrowyoffset = \arrowcenteryoffset + \arrowyoffset;
    \bottomarrowyoffset = \arrowcenteryoffset - \arrowyoffset;
    \arrowtoboxpadding = 5;
}
\coordinate (center) at (0,0);
\node (dealer) at (center) [
    draw,
    fill=white!80!gray,
    very thick,
    minimum width=3cm,
    minimum height=2cm,
    align=center,
    rounded corners=.50cm,
] {JP Morgan};

\node (counterparty) at ($(dealer.west) + (-5, 0)$) [
    draw,
    very thick,
    minimum width=3cm,
    minimum height=2cm,
    align=center
] {Counterparty};

\node (exchange) at ($(dealer.east) + (5, 0)$) [
    draw,
    dashed,
    very thick,
    minimum width=3cm,
    minimum height=2cm,
    align=center
] {Exchange};

\node (funding) at ($(dealer.south) + (0, -4)$) [
    draw,
    very thick,
    minimum width=3cm,
    minimum height=2cm,
    align=center
] {Funding desk};

% Dealer - Counterparty relations
\draw[->, thick] 
    ([yshift=\bottomarrowyoffset pt, xshift=-\arrowtoboxpadding pt]dealer.south west) -- 
    ([yshift=\bottomarrowyoffset pt, xshift= \arrowtoboxpadding pt]counterparty.south east)
    node[pos = 0.35, fill=white, inner sep=2pt] {\footnotesize
        floating leg
    };
\draw[<-, thick] 
    ([yshift=\toparrowyoffset pt, xshift=-\arrowtoboxpadding pt]dealer.south west) -- 
    ([yshift=\toparrowyoffset pt, xshift= \arrowtoboxpadding pt]counterparty.south east)
    node[pos = 0.65, fill=white, inner sep=2pt] {\footnotesize
        fixed leg
    };

% Dealer - Exchange relations
\draw[<-, thick] 
    ([yshift=\bottomarrowyoffset, xshift= \arrowtoboxpadding pt]dealer.south east) -- 
    ([yshift=\bottomarrowyoffset, xshift=-\arrowtoboxpadding pt]exchange.south west)
    node[pos = 0.65, fill=white, inner sep=2pt] {\footnotesize
        floating leg
    };
\draw[->, thick] 
    ([yshift=\toparrowyoffset, xshift= \arrowtoboxpadding pt]dealer.south east) -- 
    ([yshift=\toparrowyoffset, xshift=-\arrowtoboxpadding pt]exchange.south west)
    node[pos = 0.35, fill=white, inner sep=2pt] {\footnotesize
        fixed leg
    };

\draw[->, thick, draw=funding-color] 
    ([yshift=-\bottomarrowyoffset, xshift= \arrowtoboxpadding pt]dealer.north east) -- 
    ([yshift=-\bottomarrowyoffset, xshift=-\arrowtoboxpadding pt]exchange.north west)
    node [pos=0.4, fill=white, inner sep=0pt] {
        \includegraphics[width=0.5cm]{\graphicsfolder/source-of-funding-costs/collateral}
    };
\draw[<-, thick, draw=rate-color] 
    ([yshift=-\toparrowyoffset, xshift= \arrowtoboxpadding pt]dealer.north east) -- 
    ([yshift=-\toparrowyoffset, xshift=-\arrowtoboxpadding pt]exchange.north west)
    node [pos=0.6, fill=white, inner xsep=2pt, inner ysep=0] {
        \footnotesize OIS
    };

% Dealer - Funding desk relations
\draw[<-, thick, draw=funding-color]
    ([xshift=-\arrowyoffset, yshift=-\arrowtoboxpadding]dealer.south) --
    ([xshift=-\arrowyoffset, yshift= \arrowtoboxpadding]funding.north)
    node[midway, fill=white, inner xsep=2pt, inner ysep=0, align=center, rotate=90] {
        \footnotesize Funding%
    };
\draw[->, thick, draw=rate-color]
    ([xshift=\arrowyoffset, yshift=-\arrowtoboxpadding]dealer.south) --
    ([xshift=\arrowyoffset, yshift= \arrowtoboxpadding]funding.north)
    node[midway, fill=white, inner xsep=2pt, inner ysep=0, align=center, rotate=90] {
        \footnotesize OIS + S%
    };
            \end{tikzpicture}
            }   
        \end{figure}

        \begin{tikzpicture}
            [
                remember picture,
                overlay,
                shift={([yshift=2.2cm,xshift=3cm]current page.south west)},
                align=left
            ]
            \node[font=\small]{
                \textbf{The rate effect} \\
                
                floating leg $\downarrow$

                \\[0.2cm]

                \hspace{0.8em} \rotatebox[origin=c]{180}{$\Lsh$} \includegraphics[width=0.25cm]{\graphicsfolder/source-of-funding-costs/collateral} $\uparrow$

                \\[0.2cm]

                \hspace{1.6em} \rotatebox[origin=c]{180}{$\Lsh$} funding $\uparrow$
            };
        \end{tikzpicture}

        \begin{tikzpicture}
            [
                remember picture,
                overlay,
                shift={([yshift=2.2cm,xshift=-3cm]current page.south east)},
                align=left
            ]
            \node[font=\small]{
                \textbf{The funding effect} \\
                
                default intensity $\uparrow$

                \\[0.2cm]

                \hspace{0.8em} \rotatebox[origin=c]{180}{$\Lsh$} S $\uparrow$

                \\[0.2cm]

                \hspace{1.6em} \rotatebox[origin=c]{180}{$\Lsh$} funding rate $\uparrow$
            };
        \end{tikzpicture}

    \end{frame}
\end{document}
